\chapter{结论}
在三页课件、小班并发、预置课程数据的演示范围内，LingoBridge 平台具备如下可行性条件：
\begin{itemize}
  \item \textbf{技术方面}：通过 React 前端框架与 Node/Express 后端服务的轻量级架构，辅以自适应状态重连机制，能够支撑跨海小班课堂的基本同步及录制需求，具备基本技术开发条件；
  \item \textbf{经济方面}：基于边缘静态托管（Vercel）与轻量级计算节点（CVM），配合双层语音缓存策略，在演示及试点阶段能够有效控制基础设施和 TTS 重复生成的费用支出；
  \item \textbf{操作方面}：教师端排课及课件上传流程、学生端课堂页面同步及麦克风录音提交流程清晰，且系统预留了后续接入外部智能 ASR 评分和多语言翻译的扩展接口。
\end{itemize}

综上所述，LingoBridge 平台在上述演示与试点条件限制下，\textbf{具备有条件的可行性}。
